% Story spine for the main text (working draft, 14 August 2026):
% Similar overt choices can be generated by different internal learning
% processes. We first establish this behavioral underdetermination, then use
% withheld oral reports to adjudicate between behaviorally equivalent models.
% Once the latent rule distribution has independent support, we use it to
% explain feedback-gated rule revision, its reorganization by task structure,
% and the autonomous generation of individual learning phenotypes.
%
% The main claim would be falsified if (i) a simpler model matched both choices
% and oral reports, (ii) latent transitions failed to predict independently
% observed report or behavioral transitions, or (iii) the selected model failed
% to reproduce learning when it controlled its own choices and feedback.
%
% Magenta text marks evidence that has not yet been generated. It is deliberately
% attached to the claim it must support, rather than collected as a checklist.

\providecommand{\resultpending}[1]{\textcolor{magenta}{[\textit{Result pending: #1}]}}

\section{Results}\label{sec:results}

Our central question was not whether a sufficiently flexible model could mimic
human choices. It was whether the same choices were sufficient to identify how
a learner searched, abandoned and eventually committed to a rule. Category
learning makes this inverse problem especially acute: a correct response can
reflect knowledge of the correct partition, a partially correct rule, or a
fortuitous choice under an entirely different rule. We therefore designed the
analysis to expose this ambiguity rather than conceal it. We first identified
models that made comparably accurate choice predictions but implied different
latent learning histories. We then asked whether oral reports---withheld from
model fitting---could distinguish those histories. Only after this independent
test did we use the inferred states to explain learning dynamics and ask whether
the same process could generate them autonomously.

% Planned Fig. 1: task manipulation, diverse behavioral trajectories, oral-rule
% episodes, and examples of choice-equivalent but internally distinct trials.

\subsection{Similar learning outcomes concealed distinct paths through rule space}

Ninety-six participants learned deterministic category structures built from
the same four-dimensional stimuli. The three between-participant conditions
separated elementary rule discovery from structured generalization: condition
1 required a binary division along one feature, whereas conditions 2 and 3
required a four-category hierarchical solution whose organization was either
disclosed or left to be discovered. This design allowed endpoint performance
to be separated from the route by which that endpoint was reached.

Participants did not approach the solution as uniformly noisy versions of one
ideal learner. Aggregate accuracy improved in every condition
\resultpending{report condition-wise learning effects, uncertainty and
between-participant dispersion}, but individual trajectories contained extended
plateaux, abrupt gains, reversals and persistent misconceptions.
\resultpending{quantify change points, dwell times and the prevalence of abrupt
versus gradual learning phenotypes.} Importantly, similar accuracy curves could
mask different paths: some participants repeatedly tested alternatives before
a sudden transition, whereas others refined a stable but initially imprecise
rule. Thus, the behavioral variation to be explained concerned the organization
of learning, not merely its average rate.

Trialwise oral reports made that organization partly observable. Participants
described structured relations---including feature magnitude, ordering,
near-equality and central values---and often maintained a reported rule across
several decisions before revising it. In the binary task, adding the recurrent
near-equality and central-value relations to the initial rule catalogue yielded
29 candidate rules and increased coverage of structured reports from 84.5\% to
92.8\%. The oral data therefore supported a discrete but graded description of
what participants considered: cognition moved among interpretable regions of
rule space rather than along a single scalar measure of certainty.

The oral reports also contained information not recoverable from the current
choice alone. Among trials matched on stimulus, choice, feedback history and
current accuracy, the reported rule predicted subsequent rule revision and
choice change \resultpending{report conditional prediction or matched-trial
effect with participant-level uncertainty}. This result is important for the
logic that follows. The reports were neither treated as infallible access to
mental state nor reduced to post-hoc explanations of a response; they provided
a noisy second measurement of the evolving rule representation.

\subsection{Behaviorally successful models disagreed about what learners believed}

We next turned the inverse problem into a direct model comparison. The model
set spanned two qualitatively different accounts. Continuous-accumulation
models maintained and updated evidence across the available rule space. The
dynamic rule model instead represented a capacity-limited active set, separated
rules that were currently considered from the rule actually used for choice,
and allowed feedback-dependent failure and mastery signals to control local
refinement, global search and commitment. Nested variants removed these
operations in turn, permitting the contribution of adaptive search, persistent
execution and rule-specific confidence to be tested without changing the task
representation.

All hyperparameters were selected from choices alone by coordinate-descent
search (Hyper-CD) on the training portion of each participant's sequence. The
model structures and hyperparameters were frozen before forward prediction.
Rather than declaring a winner from a small likelihood difference, we first
defined a behavioral-equivalence set using held-out choice likelihood together
with participant-level posterior-predictive checks of accuracy, volatility,
switching and history dependence.
\resultpending{report the temporal split, Hyper-CD convergence, held-out model
contrasts and the models retained in the behavioral-equivalence set.}

Several models accounted for overt choices to a statistically and practically
similar degree \resultpending{report paired NLL differences, equivalence bounds
and posterior-predictive coverage}, yet assigned sharply different probability
to the rules a participant was considering on the same trial
\resultpending{report cross-model divergence of latent rule distributions}.
In particular, a model could reproduce when a participant was correct while
misidentifying why. Choice prediction therefore constrained the possible
learning processes but did not identify one of them.

This ambiguity was not resolved by selecting the single simulated trajectory
that most resembled a participant. For every candidate model, particle-filtered
(PF) inference approximated the distribution over all latent rule paths
compatible with the preceding choices. Lines and uncertainty bands in the
following analyses summarize this marginal distribution. The inferential object
is consequently a trialwise distribution over possible internal states, not a
retrospectively chosen narrative.

% Planned Fig. 2: held-out behavioral model comparison, PPCs, and pairs of
% behaviorally equivalent models with divergent latent rule distributions.

\subsection{Withheld oral reports broke the behavioral equivalence}

Behavioral equivalence created the decisive test of the model. Oral reports had
not contributed to participant-level parameter estimation, PF weighting or
model selection. We mapped each report to a likelihood over candidate rules and
retained the latest report about each category, updating the reported category
while preserving the other category's most recent value. Combining these
category-specific likelihoods produced an oral-report distribution over the
same rule space as the model. This representation avoided forcing an ambiguous
statement into a single rule and allowed partial information to accumulate
across trials.

We compared model and oral distributions at three levels. Full-distribution
similarity tested whether probability mass was allocated to the same competing
rules. Expected category centres tested agreement in the represented partition
without requiring identical labels. Target-based analyses summed over all rules
that implemented the reported category relation, rather than crediting only one
arbitrarily designated target rule. Time-preserving permutations and held-out
participants established how much alignment could arise from common progress
through the task or from using an oral-informed rule vocabulary.

The dynamic rule model aligned more closely with the withheld oral distributions
than did the behaviorally equivalent continuous and reduced models
\resultpending{report participant-level differences in full-distribution,
expected-centre and target-equivalence alignment, with temporal-null and
held-out-catalogue tests}. Crucially, this advantage remained after matching
models on held-out choice prediction \resultpending{report the conditional or
matched-model analysis}. The evidence thus formed a double dissociation: overt
behavior alone admitted multiple computational explanations, whereas the
independent measurement of represented rules selected among them.

The currently completed fixed-configuration analysis establishes the scale of
the measurable correspondence, although it is not a substitute for the
Hyper-CD comparison. Across condition-1 participants, expected-centre similarity
between the model and oral distributions averaged 0.885, and active-space
Jensen--Shannon similarity averaged 0.478. Trialwise target-equivalence
trajectories were positively associated \resultpending{replace the current
single-target diagnostic with the final equivalence-class estimate and its
participant-level interval}. These descriptive values show that the two
channels converged, but the stronger mechanistic conclusion depends on the
behavior-matched model contrast above.

The oral measure was treated as external evidence, not as ground truth. Its
validity depended on three separable observations: reports were semantically
related to the stimulus, stable when participants maintained a strategy and
sensitive when behavior changed.
\resultpending{report these measurement-validity analyses and robustness to
report omissions, parser uncertainty and category-specific carry-forward.}
This distinction makes the conclusion narrower but stronger: the selected
model recovered the component of latent rule use that was shared by two noisy
and independently analysed measurements.

% Planned Fig. 3: the decisive double dissociation---similar choice fit but
% different oral alignment---plus marginal PF and oral distributions over time.

\subsection{Rule revision was punctuated and gated by the consequences of failure}

Having established which latent distribution was psychologically constrained,
we asked what it revealed about the learning algorithm. The dynamic model made
a temporal prediction that a smooth evidence accumulator did not. Negative
feedback should not cause an undifferentiated loss of confidence. Initially it
should increase failure pressure within the current neighbourhood, prompting
local refinement; persistent failure should then reorganize the active set by
recruiting distant rules. Only after this representational change should the
executed rule, and therefore overt choice policy, switch. Conversely, sustained
success should increase mastery, stabilize execution and suppress global
search.

Event-locked analyses supported this sequence. Failure pressure rose before
active-set replacement, active-set replacement preceded executed-rule switches,
and those switches anticipated changes in oral reports and behavior
\resultpending{report event-locked time courses, lead--lag contrasts and
participant-level uncertainty}. Local substitutions dominated after isolated
errors, whereas error streaks selectively increased global transitions
\resultpending{report transition-distance effects and their interaction with
error-streak length}. Following successful evidence, the probability of
retaining the executed rule increased and rule uncertainty contracted
\resultpending{report mastery and dwell-time effects}.

These results distinguish an algorithmic account from a relabelling of the
learning curve. The relevant model variables predicted independently measured
transitions at their specified temporal positions, and reduced models lacking
the corresponding control operation failed selectively at those events
\resultpending{report event-specific ablations}. Human learning was therefore
better characterized as punctuated, feedback-gated reorganization of a bounded
rule workspace than as uniformly paced convergence on a fixed representation.

% Planned Fig. 4: event-locked chain from feedback to failure/mastery, active-set
% reorganization, executed-rule switch, oral switch and behavioral consequence.

\subsection{Task structure reorganized search rather than merely slowing learning}

The hierarchy manipulation provided a test of whether the inferred mechanism
generalized beyond the binary task. We kept the learning architecture fixed and
changed only the rule library and information made available by the task. If
the model merely absorbed idiosyncratic condition-1 choices, it should fail to
predict how search was reorganized when the solution required a hierarchy.

Disclosing the hierarchy constrained global search and promoted staged local
refinement within the relevant branch, whereas hiding the same hierarchy
prolonged uncertainty over branches and increased returns to global search
\resultpending{report cross-condition effects on search distance, active-set
entropy, branch errors and executed-rule dwell time}. These differences
explained condition effects on learning time beyond perceptual discriminability
and mean response noise \resultpending{report mediation or model-based
decomposition, using the independent perceptual calibration}. The same
feedback-gated control law thus adapted its search to the structure of the
problem rather than requiring a different learning algorithm for each task.

This comparison also tests a resource-rational interpretation of the dynamics.
Participants facing greater independently measured uncertainty should require
more evidence before commitment and should return to global search more often,
but the direction and timing of these effects must follow from recoverable
model quantities rather than from unconstrained parameter correlations.
\resultpending{report the preregistered relation to perceptual uncertainty and
rule out learning-duration and task-condition confounds.}

\subsection{The validated process generated individual learning phenotypes autonomously}

The preceding analyses conditioned inference on human choices. They could show
that the model tracked latent states, but not that its learning process was
sufficient to create a human-like learner. We therefore closed the explanatory
loop with autonomous simulation. On each trial, the fitted model sampled its
own category choice, received the feedback caused by that choice and updated
its later state from its own history. No human choice was reintroduced after
initialization.

The selected model reproduced not only mean accuracy but the joint distribution
of time to criterion, abrupt and gradual improvement, misconception dwell,
local volatility and rule-switch timing across participants
\resultpending{report autonomous posterior-predictive coverage and uncertainty
for each preregistered phenotype}. It also generated the observed coupling
between oral-rule transitions and behavioral change when synthetic reports were
sampled from the independently validated observation model
\resultpending{report this out-of-sample generative check}. Reduced models made
diagnostic errors: continuous accumulation converged too uniformly, removal of
adaptive search eliminated prolonged misconceptions or recovery from them, and
removal of persistent execution produced excessive trial-to-trial switching
\resultpending{report the model-by-phenotype failure matrix}.

Conditional trajectory replay is not this generative test. Replay continues to
ingest a participant's choices and feedback and asks how latent paths would
evolve under that observed history; autonomous simulation lets the model's own
choices alter every subsequent state. The former is useful for diagnosing
inference and appears only as a control, whereas the latter supports the claim
that the proposed process can produce the behavioral phenomena it explains.

% Planned Fig. 5: autonomous posterior-predictive phenotypes and selective
% failures of models that were difficult to separate with one-step choice fit.

\subsection{Recovery analyses bounded the level of mechanistic interpretation}

Finally, we asked which parts of this account were identifiable from sequences
of the present length. Structure recovery was evaluated before parameter
recovery: data generated by each finalist model were passed through the entire
Hyper-CD and model-selection pipeline, testing whether the architecture could
be distinguished rather than whether one parameter could be recovered while
the correct architecture was assumed.
\resultpending{report the cross-generating-model confusion matrix and the range
of conditions in which the full structure is recoverable.}

The targeted recovery analysis already completed for three control parameters
showed why this distinction matters. Across 72 autonomously generated
condition-1 trajectories and nine joint parameter profiles, the exact profile
was recovered in 13 cases (18.1\%; 95\% Wilson interval, 10.9--28.5\%), compared
with a chance rate of 11.1\%. Individual-parameter recovery ranged from 37.5\%
to 43.1\%, with every interval including the one-in-three chance rate. We
therefore do not interpret these fitted values as precise participant-level
psychological traits. Parameter-level claims will be retained only where the
final design demonstrates recovery; otherwise inference remains at the level
of recoverable model structure, latent functions and event-level predictions.

PF likelihood calibration provides a second boundary. Candidate rankings were
not stable under the completed 32--256-particle audit, so the formal Hyper-CD
and likelihood-ranking results above remain placeholders until a stable
particle budget or lower-variance estimator is established. This numerical
control concerns confidence in model selection, not the cognitive distinction
between a distribution over latent paths and a selected best trajectory.

Taken together, the intended evidence chain goes beyond showing that a model
fits learning curves. Choices first reveal the underdetermination of the hidden
learning process; a withheld oral channel then adjudicates among the competing
processes; independently observed transitions test the proposed control law;
task structure tests its generality; and autonomous simulation tests its
generative sufficiency. On this account, individual category learning is not
simply noisy convergence toward a known answer. It is the feedback-gated
construction, testing and stabilization of a rule within a bounded cognitive
workspace.
